\section{OBJETIVOS ESPECÍFICOS}

\begin{itemize}
\item Investigar las operaciones de mecanizado en torno CNC, los mecanismos de desgaste de herramientas de corte y las tecnologías de inteligencia artificial en el borde (\textit{Edge AI}) orientadas al monitoreo de condición.
\item Diseñar la arquitectura mecatrónica del sistema periférico de monitoreo, integrando acondicionamiento de señal, procesamiento en el borde y conectividad IoT.
\item Dimensionar los sensores de corriente, vibración y la unidad de procesamiento embebido.
\item Integrar el sistema periférico de adquisición, procesamiento inteligente y supervisión remota.
\item Validar experimentalmente el funcionamiento del sistema integrado mediante pruebas de corte en los tornos CNC de MAINSA.
\end{itemize}

\subsection{Acciones del proyecto}
Con el propósito de operativizar cada uno de los objetivos específicos y detallar la metodología de ejecución técnica, en la \tab{acciones_objetivos} se desglosan las acciones concretas asociadas a cada fase del proyecto.

{\setlength{\parskip}{1.5pt}
\begin{longtable}{|>{\raggedright\arraybackslash}p{4.8cm}|>{\raggedright\arraybackslash}p{9.2cm}|}
\caption{Acciones asociadas a los objetivos específicos\label{tab:acciones_objetivos}}
\\
\hline
\textbf{Objetivo específico} & \textbf{Acciones a desarrollar} \\
\hline \hline
\endfirsthead

\hline
\textbf{Objetivo específico} & \textbf{Acciones a desarrollar} \\
\hline \hline
\endhead

\hline
\endfoot

\multicolumn{2}{c}{\textbf{Fuente:} Elaboración propia}
\endlastfoot

Investigar operaciones de mecanizado en torno CNC, fenomenología del desgaste de herramientas y tecnologías de Edge AI para monitoreo. & 
• Revisión y caracterización de las operaciones de torneado CNC (cilindrado, refrentado, ranurado) y cinemática de corte.\par\vspace{1mm}
• Análisis de la fenomenología del desgaste de herramientas y patrones de falla según la norma ISO 3685.\par\vspace{1mm}
• Estudio del estado del arte en sistemas de monitoreo indirecto de herramientas (TCM) mediante corriente y vibración.\par\vspace{1mm}
• Evaluación de tecnologías de Inteligencia Artificial en el borde (\textit{Edge AI} / TinyML) y plataformas embebidas.\\ \hline

Diseñar la arquitectura mecatrónica del sistema periférico de monitoreo, acondicionamiento, Edge AI e IoT. & 
• Definición del diagrama de bloques funcional (adquisición, acondicionamiento, procesamiento y enlace IoT).\par\vspace{1mm}
• Diseño del circuito de acondicionamiento para rectificación y conversión a niveles proporcionales al valor eficaz (RMS).\par\vspace{1mm}
• Diseño esquemático y ruteo de la placa de circuito impreso (PCB) con aislamiento galvánico industrial.\par\vspace{1mm}
• Concepción del gabinete protector y soportes de fijación no invasiva para la instalación en la máquina.\\ \hline

Dimensionar los sensores de corriente, vibración y la unidad de procesamiento embebido. & 
• Dimensionamiento del rango de medición y relación de transformación del transductor de corriente según la potencia del motor.\par\vspace{1mm}
• Selección del acelerómetro considerando sensibilidad, ancho de banda y rango dinámico de vibración en la torreta.\par\vspace{1mm}
• Selección del microcontrolador evaluando periféricos ADC, DMA, instrucciones DSP y memoria requerida.\par\vspace{1mm}
• Cálculo y selección de componentes pasivos de filtrado analógico, resistencias de carga y fuentes de alimentación.\\ \hline

Integrar el sistema periférico de adquisición, procesamiento inteligente y supervisión remota. & 
• Ensamblaje, integración eléctrica y verificación funcional de la placa de circuito impreso y módulos acondicionadores.\par\vspace{1mm}
• Programación del firmware para muestreo síncrono, filtrado digital y extracción de descriptores estadísticos temporales.\par\vspace{1mm}
• Entrenamiento, cuantización y despliegue del modelo de aprendizaje automático ligero para clasificación de desgaste.\par\vspace{1mm}
• Configuración del protocolo de comunicaciones y desarrollo del panel de supervisión remota para visualización y alertas.\\ \hline

Validar experimentalmente el funcionamiento del sistema integrado mediante pruebas de corte en MAINSA. & 
• Instalación y calibración del sistema periférico en el torno CNC bajo condiciones operativas de producción real.\par\vspace{1mm}
• Ensayos de torneado cilíndrico registrando señales con insertos en estado nuevo, desgaste moderado y crítico.\par\vspace{1mm}
• Medición metrológica de la rugosidad superficial ($Ra$) en probetas y verificación de tolerancias dimensionales.\par\vspace{1mm}
• Evaluación cuantitativa de la exactitud de clasificación y cuantificación de la reducción de paradas y mermas.\\ \hline

\end{longtable}
}
